\section*{Capítulo \ref{ch:b1:isostasy-gravity} --- Isostasia e anomalias da gravidade}

\begin{solution}{exo:b1:isostasy-gravity:1}
(a) $\qty{-150}{\milli\gal} = -150\times\qty{e-5}{m.s^{-2}} = \qty{-1.5e-3}{m.s^{-2}}$,
e $\num{1.5e-3}/9.81 = \num{1.53e-4}$ de $g$. (b) $\qty{1}{\milli\gal} =
\qty{e-5}{m.s^{-2}}$, \num{1.02e-6} de $g$; $\qty{1}{\micro\gal} =
\qty{e-8}{m.s^{-2}}$, \num{1.02e-9} de $g$: um gravímetro que resolve um microgal mede a gravidade com cerca de uma parte em um bilhão.
\end{solution}

\begin{solution}{exo:b1:isostasy-gravity:2}
(a) No equador ambos os senos se anulam: $\gamma = \qty{9.780327}{m.s^{-2}}$. A
\ang{45}, $\sin^2\varphi = 0.5$ e $\sin^2 2\varphi = 1$:
$\gamma = 9.780327\times(1 + 0.0026512 - 0.0000058) = \qty{9.806200}{m.s^{-2}}$.
Nos polos, $\sin^2\varphi = 1$ e $\sin 2\varphi = 0$:
$\gamma = 9.780327\times1.0053024 = \qty{9.832186}{m.s^{-2}}$.
(b) $9.832186 - 9.780327 = \qty{0.051859}{m.s^{-2}}$, isto é
\qty{5186}{\milli\gal}. (c) $3392/5186 = 0.65$: o efeito centrífugo explica
cerca de dois terços da diferença; o achatamento, que aproxima os polos do
centro, o resto.
\end{solution}

\begin{solution}{exo:b1:isostasy-gravity:3}
$r = h\rho_c/(\rho_m - \rho_c) = 3\times2800/500 = \qty{16.8}{km}$. A \omterm{def:b1:layered-earth:layers}{crosta} é
mais espessa em $h + r = \qty{19.8}{km}$ (isto é $h\rho_m/(\rho_m - \rho_c) =
3\times3300/500$).
\end{solution}

\begin{solution}{exo:b1:isostasy-gravity:4}
(a) $1/0.3086 = \qty{3.24}{m}$. (b) $0.01/0.3086 = \qty{0.032}{m}$: as alturas
devem ser conhecidas a cerca de \qty{3}{cm}, por isso um levantamento registra a altura de cada
estação com posicionamento por satélite ou nivelamento.
\end{solution}

\begin{solution}{exo:b1:isostasy-gravity:5}
(a) $\rho = 2700\times100/(100 + h)$: \qty{2673}{kg.m^{-3}} para \qty{1}{km},
\qty{2621}{kg.m^{-3}} para \qty{3}{km}, \qty{2571}{kg.m^{-3}} para \qty{5}{km}.
(b) $\rho = (2700\times100 - 1030\times4)/96 = \num{265880}/96 =
\qty{2770}{kg.m^{-3}}$: a coluna sob o mar é mais densa. (c) O modelo de Pratt
pede variações laterais de densidade a profundidade constante. Sob o fundo oceânico
essas variações existem: \omterm{def:b1:layered-earth:lithosphere}{litosfera} jovem e quente perto de uma dorsal é mais leve que \omterm{def:b1:layered-earth:lithosphere}{litosfera}
velha e fria longe dela (\cref{ch:b1:heat-budget}). Sob os continentes, ao contrário,
a espessura da \omterm{def:b1:layered-earth:layers}{crosta} se mostra variar com a topografia.
\end{solution}

\begin{solution}{exo:b1:isostasy-gravity:6}
(a) Estação 1: $g - \gamma = 9.807751 - 9.808911 = \qty{-0.001160}{m.s^{-2}} =
\qty{-116.0}{\milli\gal}$; \omterm{def:b1:isostasy-gravity:anomalies}{correção do ar livre} $0.3086\times350 =
\qty{108.0}{\milli\gal}$, logo $\Delta g_{\mathrm{FA}} = \qty{-8.0}{\milli\gal}$;
\omterm{def:b1:isostasy-gravity:anomalies}{correção de Bouguer} $0.112\times350 = \qty{39.2}{\milli\gal}$, logo
$\Delta g_{\mathrm B} = \qty{-47.2}{\milli\gal}$.
Estação 2: $g - \gamma = \qty{-535.9}{\milli\gal}$; \omterm{def:b1:isostasy-gravity:anomalies}{correção do ar livre}
$0.3086\times1850 = \qty{570.9}{\milli\gal}$, logo $\Delta g_{\mathrm{FA}} =
\qty{+35.0}{\milli\gal}$; \omterm{def:b1:isostasy-gravity:anomalies}{correção de Bouguer} $0.112\times1850 =
\qty{207.2}{\milli\gal}$, logo $\Delta g_{\mathrm B} = \qty{-172.2}{\milli\gal}$.
(b) Na estação 2 a \omterm{def:b1:isostasy-gravity:anomalies}{anomalia de Bouguer} é fortemente negativa: sob as montanhas há um
déficit de massa, uma raiz de \omterm{def:b1:layered-earth:layers}{crosta} leve. A \omterm{def:b1:isostasy-gravity:anomalies}{anomalia do ar livre}, pequena em comparação com a
\omterm{def:b1:isostasy-gravity:anomalies}{correção de Bouguer}, mostra que o déficit quase equilibra as montanhas: elas estão em grande parte compensadas
(cerca de $172/207$, ou \qty{83}{\percent}, se o modelo de placa valer para toda a cadeia).
\end{solution}

\begin{solution}{exo:b1:isostasy-gravity:7}
(a) Como na demonstração da proposição, uma camada fina à distância $z$
contribui com a atração $G\rho\,2\pi z\,h\int_0^a s\,(s^2 + z^2)^{-3/2}\,\mathrm ds$, agora
com a integral parada em $s = a$:
$\int_0^a s\,(s^2 + z^2)^{-3/2}\,\mathrm ds = 1/z - 1/\sqrt{z^2 + a^2}$. Logo
$g = 2\pi G\rho h\,\bigl(1 - z/\sqrt{z^2 + a^2}\bigr)$, que tende ao valor da placa
quando $a \to \infty$. (b) $z/\sqrt{z^2 + a^2} = 0.1$ dá
$\sqrt{z^2 + a^2} = 10z$, $a = \sqrt{99}\,z = \qty{9.95}{km}$ para
$z = \qty{1}{km}$. (c) A rocha dentro de cerca de dez vezes sua profundidade abaixo da
estação fornece \qty{90}{\percent} da atração: a correção de terreno
precisa da topografia em um raio de vários quilômetros a dezenas de quilômetros em torno de cada
estação, não apenas embaixo dela.
\end{solution}

\begin{solution}{exo:b1:isostasy-gravity:8}
(a) $r' = d(\rho_c - \rho_w)/(\rho_m - \rho_c) = 4\times1770/500 =
\qty{14.16}{km}$, e a \omterm{def:b1:layered-earth:layers}{crosta} tem $35 - 4 - 14.16 = \qty{16.84}{km}$ de espessura.
(b) O modelo de Airy com uma só densidade crustal prevê uma \omterm{def:b1:layered-earth:layers}{crosta} oceânica espessa demais em mais de
três vezes espessa demais. A \omterm{def:b1:layered-earth:layers}{crosta} oceânica não é uma \omterm{def:b1:layered-earth:layers}{crosta} continental afinada: é
uma rocha diferente, mais densa, e os dois níveis da Terra são dois tipos de
\omterm{def:b1:layered-earth:layers}{crosta}, como sugere a \omterm{def:b1:isostasy-gravity:hypsometry}{curva hipsométrica}.
\end{solution}

\begin{solution}{exo:b1:isostasy-gravity:9}
(a) Um segundo de arco é $\pi/(180\times3600) = 1/\num{206265}$ de radiano:
\qty{5.236}{''} $= \qty{2.538e-5}{rad}$ e \qty{15.885}{''} $=
\qty{7.701e-5}{rad}$. (b) $a_h = g\theta$: $9.79\times\num{2.538e-5} =
\qty{2.485e-4}{m.s^{-2}} = \qty{24.9}{\milli\gal}$ observados, e
$9.79\times\num{7.701e-5} = \qty{75.4}{\milli\gal}$ calculados a partir das
montanhas visíveis. (c) $1 - 5.236/15.885 = 0.670$: a compensação cancelou cerca de dois
terços da atração das montanhas visíveis.
\end{solution}

\begin{solution}{exo:b1:isostasy-gravity:10}
(a) Fora dela, uma esfera atrai como uma massa pontual em seu centro, com
$GM/(x^2 + z^2)$ ao longo da linha até o centro; a componente vertical é isso
vezes o cosseno $z/\sqrt{x^2 + z^2}$, o que dá $GMz/(x^2 + z^2)^{3/2}$.
(b) $M = \tfrac43\pi(\qty{2000}{m})^3\times(\qty{-200}{kg.m^{-3}}) =
\qty{-6.70e12}{kg}$, e acima do centro $GM/z^2 = \num{6.674e-11}\times
(\num{-6.70e12})/(4000)^2 = \qty{-2.80e-5}{m.s^{-2}} = \qty{-2.80}{\milli\gal}$.
(c) $(1 + x^2/z^2)^{-3/2} = 1/2$ dá $x = z\sqrt{2^{2/3} - 1} = 0.766\,z =
\qty{3.07}{km}$. (d) A semilargura à meia altura do perfil dá a
profundidade: $z = 1.305\,x_{1/2}$. Um corpo mais profundo dá uma anomalia mais larga e mais suave.
\end{solution}

\begin{solution}{exo:b1:isostasy-gravity:11}
(a) $V = \tfrac13\pi R^2 H = \tfrac13\pi(\qty{15e3}{m})^2(\qty{3e3}{m}) =
\qty{7.07e11}{m^3}$, e $\Delta m = (2800 - 1030)V = \qty{1.25e15}{kg}$.
(b) O centro de massa fica \qty{0.75}{km} acima do fundo, a
$d = 5 - 0.75 = \qty{4.25}{km}$ abaixo da superfície do mar. Ali a anomalia de potencial
é $T = -G\Delta m/d$, logo $N = G\Delta m/(\gamma d) =
\num{6.674e-11}\times\num{1.25e15}/(9.81\times4250) = \qty{2.0}{m}$.
(c) A maior parte da massa do cone está em sua parte inferior larga, a 5 a
\qty{16}{km} do ponto acima do cume: a média de $1/r$ sobre a massa é menor que $1/\qty{4.25}{km}$, e a massa pontual superestima o
potencial. Somando discos finos obtém-se cerca de \qty{1.1}{m}.
(d) Uma raiz de \omterm{def:b1:layered-earth:layers}{crosta} leve embaixo poria um déficit da mesma massa mais fundo;
os dois potenciais quase se cancelam, e a elevação seria muito menor,
mas não nula, já que o déficit fica mais longe da superfície que o excesso.
\end{solution}

\begin{solution}{exo:b1:isostasy-gravity:12}
(a) A taxa de soerguimento é $-\mathrm dw/\mathrm dt = w_0\,e^{-t/\tau}/\tau = w/\tau$.
(b) $w = 818\,e^{-10000/4000} = 818\times0.0821 = \qty{67.1}{m}$, e a taxa é
$\qty{67.1}{m}/\num{4000}\,\unit{yr} = \qty{16.8}{mm.yr^{-1}}$, mais de uma
vez e meia o valor medido. (c) Com $\tau = 3100$ anos, $w =
818\,e^{-3.226} = \qty{32.5}{m}$ e a taxa é \qty{10.5}{mm.yr^{-1}}; com
$\tau = \num{68000}$ anos, $w = 818\,e^{-0.147} = \qty{706}{m}$ e a taxa
é \qty{10.4}{mm.yr^{-1}}. As duas implicam depressões remanescentes muito diferentes.
Uma depressão não compensada é um déficit de massa de rocha do \omterm{def:b1:layered-earth:layers}{manto}: sobre uma região ampla ela dá uma \omterm{def:b1:isostasy-gravity:anomalies}{anomalia do ar livre} de cerca de $-2\pi G\rho_m w$, isto é
$-0.0419\times3.3\times706 \approx \qty{-98}{\milli\gal}$ para o caso lento,
contra \qty{-4}{\milli\gal} para o rápido. Um levantamento gravimétrico (ou a altura de antigas praias) as distingue. (d) A constante de tempo mede a rapidez com que o \omterm{def:b1:layered-earth:layers}{manto}
volta a fluir sob a depressão: sua viscosidade.
\end{solution}

\begin{solution}{pb:b1:isostasy-gravity:1}
\textbf{1.} $r = h\rho_c/(\rho_m - \rho_c) = 4\times2700/600 = \qty{18}{km}$.

\textbf{2.} A \omterm{def:b1:layered-earth:layers}{crosta} tem $35 + 4 + 18 = \qty{57}{km}$ de espessura; sua base, o \omterm{def:b1:layered-earth:discontinuity}{Moho},
fica a $35 + 18 = \qty{53}{km}$ abaixo do nível do mar.

\textbf{3.} Planalto: \qty{57}{km} de \omterm{def:b1:layered-earth:layers}{crosta} e $60 - 53 = \qty{7}{km}$ de \omterm{def:b1:layered-earth:layers}{manto},
$\num{57000}\times2700 + \num{7000}\times3300 = \num{1.539e8} + \num{2.31e7} =
\qty{1.770e8}{kg.m^{-2}}$. Planície: $\num{35000}\times2700 + \num{25000}\times3300
= \num{9.45e7} + \num{8.25e7} = \qty{1.770e8}{kg.m^{-2}}$. Elas são iguais, como
\cref{prop:b1:isostasy-gravity:equal-mass} exige.

\textbf{4.} $p = g\times\text{massa por unidade de área} = 9.81\times\num{1.770e8} =
\qty{1.74e9}{Pa}$, cerca de \qty{1.7}{GPa}, sob as duas colunas.

\textbf{5.} A raiz deve compensar, por sua densidade menor, a massa do
planalto: cada metro de raiz retira $\rho_m - \rho_c = \qty{600}{kg.m^{-3}}$ de
massa, e cada metro de planalto acrescenta \qty{2700}{kg.m^{-3}}, logo $2700/600 = 4.5$
metros de raiz são necessários. Com uma \omterm{def:b1:layered-earth:layers}{crosta} de \qty{2900}{kg.m^{-3}}, a razão
seria $2900/400 = 7.25$: quanto menor o contraste, mais funda a raiz.

\textbf{6.} O \omterm{def:b1:layered-earth:discontinuity}{Moho}, o forte aumento da velocidade sísmica, seria encontrado
aprofundando-se sob a borda do planalto, de \qty{35}{km} sob a planície
para cerca de \qty{53}{km} sob o planalto.

\textbf{7.} $\rho = 2700\times100/104 = \qty{2596}{kg.m^{-3}}$.

\textbf{8.} $\num{104000}\times2596.2 = \qty{2.700e8}{kg.m^{-2}}$ e
$\num{100000}\times2700 = \qty{2.700e8}{kg.m^{-2}}$.

\textbf{9.} $2700 - 2596 = \qty{104}{kg.m^{-3}}$, ou $104/2700 =
\qty{3.85}{\percent}$.

\textbf{10.} $\Delta\rho/\rho = \alpha\,\Delta T$ dá $\Delta T =
0.0385/\num{3e-5} = \qty{1280}{K}$, em média sobre toda a
coluna de \qty{104}{km}. Isso derreteria boa parte da coluna: o calor sozinho não torna este planalto leve o bastante, e o mecanismo de Pratt não pode ser o principal aqui.

\textbf{11.} As colunas de Pratt não têm raiz: o \omterm{def:b1:layered-earth:discontinuity}{Moho} permanece em cerca de \qty{35}{km}
sob o planalto, como sob a planície.

\textbf{12.} \qty{52}{km} está próximo dos \qty{53}{km} de Airy e longe dos
\qty{35}{km} de Pratt: o planalto está compensado principalmente por uma raiz crustal.

\textbf{13.} $0.0419\times2.7\times4000 = \qty{452.5}{\milli\gal}$.

\textbf{14.} Sem compensação (\cref{prop:b1:isostasy-gravity:broad}):
$\Delta g_{\mathrm{FA}} \approx \qty{+452.5}{\milli\gal}$,
$\Delta g_{\mathrm B} \approx 0$.

\textbf{15.} Compensado: $\Delta g_{\mathrm{FA}} \approx 0$ e
$\Delta g_{\mathrm B} \approx \qty{-452.5}{\milli\gal}$. A massa ausente da raiz
dá $-0.0419\times0.6\times\num{18000} = \qty{-452.5}{\milli\gal}$, a mesma,
pois $\Delta\rho\,r = \rho_c h$.

\textbf{16.} $\sin^2\ang{30} = 0.25$ e $\sin^2\ang{60} = 0.75$:
$\gamma = 9.780327\times(1 + 0.0013256 - 0.0000044) = \qty{9.793249}{m.s^{-2}}$.
Então $g - \gamma = 9.781085 - 9.793249 = \qty{-1216.4}{\milli\gal}$; a \omterm{def:b1:isostasy-gravity:anomalies}{correção do ar livre}
é $0.3086\times4000 = \qty{1234.4}{\milli\gal}$, logo
$\Delta g_{\mathrm{FA}} = \qty{+18.0}{\milli\gal}$, e
$\Delta g_{\mathrm B} = 18.0 - 452.5 = \qty{-434.5}{\milli\gal}$.

\textbf{17.} $434.5/452.5 = 0.96$: o planalto está cerca de \qty{96}{\percent}
compensado (a pequena anomalia positiva do ar livre é a parcela não compensada,
\qty{4}{\percent}).

\textbf{18.} Não. Perto do centro de um planalto amplo, cada camada anômala atua
como uma placa, cuja atração não depende de sua profundidade: só importa o déficit de massa
por unidade de área, e ele é o mesmo, $\rho_c h$, nos dois modelos. Os dois
diferem apenas perto das bordas, onde a profundidade das massas de compensação muda
a forma da anomalia.

\textbf{19.} Longe da borda, dos dois lados, a \omterm{def:b1:isostasy-gravity:anomalies}{anomalia do ar livre} é quase nula.
A alguns quilômetros para dentro da borda, o planalto sob a estação está perto e
estende-se muito para um lado apenas: ele ainda puxa com quase todo o seu valor de placa.
A raiz, a \qty{35}{km} ou mais abaixo, é vista de lá como um corpo que termina
quase em linha reta embaixo: ela puxa com pouco mais da metade de seu efeito. A
diferença é um pico positivo. A alguns quilômetros para fora, o planalto puxa
muito pouco, enquanto a raiz ainda puxa com quase metade de seu efeito: um vale negativo.
O perfil é um par positivo--negativo através de cada borda, como na figura
do capítulo.

\textbf{20.} A \omterm{def:b1:layered-earth:layers}{crosta} agora tem $57 - 1 = \qty{56}{km}$ de espessura, $\qty{21}{km}$ a mais
que a da planície. Em equilíbrio esse excesso se divide em $h' + r'$, com
$r' = 4.5\,h'$: $h' = 21/5.5 = \qty{3.818}{km}$ e $r' = \qty{17.18}{km}$.

\textbf{21.} O planalto agora se ergue a \qty{3818}{m}.

\textbf{22.} A rocha agora na superfície estava \qty{1}{km} abaixo da antiga superfície,
a \qty{3000}{m}; agora está a \qty{3818}{m}: subiu \qty{818}{m},
isto é $1000\times\rho_c/\rho_m$.

\textbf{23.} Cada quilômetro erodido abaixa a superfície em $1 - 2700/3300 =
0.182$ de quilômetro, logo baixar a superfície em \qty{4}{km} exige
$4/0.182 = \qty{22}{km}$ de erosão, toda a espessura excedente $h + r$. Em um
cinturão de montanhas antigo e desgastado, as rochas da superfície estavam, portanto, outrora a cerca de
vinte quilômetros de profundidade: as raízes profundas do cinturão estão expostas.

\textbf{24.} A superfície desce $1000 - 818$: para cada quilômetro de rocha
erodido, o abaixamento líquido da superfície é
\textbf{\qty[text-series-to-math=true]{182}{m}}.
\end{solution}
